% !TEX root = ../macro.tex
\section{Harrod--Domar: Investment and Growth}\label{sec:harrod-domar}

Investment is spending on current production and a source of future productive
capacity. Once capacity changes, a demand level sufficient to employ resources
today may be insufficient tomorrow. Harrod and Domar approach this moving
problem from different directions. We develop their questions separately before
connecting them to the fixed-proportions benchmark that Solow will change.

\subsection{From a spending equilibrium to a growth problem}

Start with a closed economy without government. At a given date, planned
investment supports demand through the multiplier. But a machine purchased
today can also produce goods later. Repeating the static spending calculation
while keeping productive capacity unchanged would discard one of investment's
two effects.

The next question is therefore about a path. How quickly must spending expand
to use the capacity being created? And would firms' responses bring the economy
toward that path if actual growth differed from what they expected? These are
questions about consistency and adjustment, before they are questions about an
optimal saving rate.

The exposition below uses continuous time, a constant saving propensity
$s\in(0,1)$, and no depreciation initially. Thus gross and net investment
coincide, $\dot K=I$. Prices are held fixed for the demand calculation. The
notation and the explicit adjustment example are teaching devices, not a
reconstruction of either author's private sequence of thought.

\begin{modelsetup}
\begin{modelcondition}{01}{Common accounting}
Output is consumed or invested. Realized saving and investment are equal:
\[
Y=C+I,\qquad S=sY=I.
\]
\end{modelcondition}
\begin{modelcondition}{02}{Domar's question}
Investment supports demand and adds capacity. What path keeps the two
consistent?
\[
Y^d=I/s,\qquad \dot Q=\kappa I.
\]
\end{modelcondition}
\begin{modelcondition}{03}{Harrod's question}
Desired investment responds to expected expansion. What growth rate validates
firms' plans?
\[
I^d=v_r\dot Y^e.
\]
\end{modelcondition}
\begin{modelcondition}{04}{Two distinct tests}
Does actual growth return to the warranted rate? Can warranted growth persist
at full employment?
\[
g_a\ \text{vs.}\ g_w,\qquad g_w\ \text{vs.}\ g_n.
\]
\end{modelcondition}
\end{modelsetup}

The two closures in the map are studied in turn. They are not imposed together
as a complete off-equilibrium model. In particular, realized investment $I$ and
desired investment $I^d$ need not coincide when expectations are disappointed.

\subsection{Domar: demand must keep up with capacity}

\subsubsection{A flow affects demand; its accumulation affects capacity}

Let $Q(t)$ be productive capacity, measured as a flow of potential real output.
Let $\kappa>0$ be the increase in capacity generated by a unit of investment.
Domar's central distinction is between the demand and capacity effects of
investment.\footnote{Evsey D. Domar (1946), ``Capital Expansion, Rate of Growth,
and Employment,'' \textit{Econometrica}, 14(2), 137--147, especially pp.~139--141.
\url{https://doi.org/10.2307/1905364}. We write $s$ for the saving propensity
and $\kappa$ for the capacity effect, rather than using his original symbols.}
The supply-side relation is
\begin{equation}\label{eq:domar-capacity}
\dot Q(t)=\kappa I(t).
\end{equation}
$\kappa$ measures added productive capacity, not the financial return earned by
an investor. With positive investment and constant $\kappa$, capacity expands
even when investment itself is not growing.

For the demand side, write desired consumption as $C^d=(1-s)Y^d$ and treat
investment as the autonomous expenditure at that date. The multiplier gives
\begin{align}
Y^d&=(1-s)Y^d+I,\\
sY^d&=I,\qquad Y^d=\frac{I}{s},\\
\dot Y^d&=\frac{\dot I}{s}.\label{eq:domar-demand}
\end{align}
This determines demand-supported output at fixed prices when capacity is
sufficient. It does not guarantee that output above $Q$ can be produced.

The distinction between $I$ and $\dot I$ is crucial. The \emph{level} of
investment adds capacity; an \emph{increase} in investment raises the
multiplier-supported level of demand. Keeping investment at a high but
constant level can therefore leave demand unchanged while capacity keeps rising.

\subsubsection{Deriving the required investment path}

Suppose demand initially matches capacity: $Y^d(0)=Q(0)$. To preserve that
equality, their changes must also match:
\begin{align}
\dot Y^d&=\dot Q,\\
\frac{\dot I}{s}&=\kappa I,\\
\frac{\dot I}{I}&=s\kappa\equiv g_D.\label{eq:domar-required-growth}
\end{align}
For positive $I(0)$, this differential equation has the solution
\begin{equation}
I(t)=I(0)e^{g_Dt},\qquad Y^d(t)=\frac{I(0)}{s}e^{g_Dt}.
\end{equation}
Integrating the capacity equation gives
\begin{equation}
Q(t)=Q(0)+\frac{I(0)}{s}\left(e^{g_Dt}-1\right).
\end{equation}
The initial condition $Q(0)=I(0)/s$ makes $Q(t)=Y^d(t)$ at every date. Matching
the derivatives alone would preserve an initial level gap, not remove it.

\begin{blueblock}
\textbf{A numerical example.}
Measure output and investment flows in units per year. Let $s=0.2$,
$\kappa=0.25$ per year, and $I(0)=20$. Initially,
\[
Y^d(0)=100,\qquad \dot Q(0)=0.25\times20=5,\qquad
\dot I(0)=0.2\times5=1.
\]
The derivatives measure changes in those annual flows per year. Investment
must initially grow at $1/20=5\%$ per year to raise demand at the same rate
as capacity. These numbers illustrate the equations; they are not estimates.
\end{blueblock}

If investment instead stays at $\bar I>0$, with $Q(0)=\bar I/s$, then
\begin{equation}
Y^d(t)=\frac{\bar I}{s},\qquad
Q(t)=\frac{\bar I}{s}+\kappa\bar I t,\qquad
\frac{Y^d(t)}{Q(t)}=\frac{1}{1+s\kappa t}.
\end{equation}
Utilization falls even though positive investment continues. A construction
expansion can support current spending while creating capacity that needs
future demand. Whether a particular investment program faces this problem
depends on its capacity effect and the demand it generates; the equations do
not establish that every program creates excess capacity.

Conversely, if demand tries to grow faster than capacity from a fully utilized
position, the multiplier calculation runs into a supply constraint. Prices,
imports, or other adjustments would require mechanisms absent here. We have
derived a path that keeps demand and capacity consistent, not a rule ensuring
that private investment follows it or a complete theory of unemployment.

\subsection{Harrod: which growth rate validates investment plans?}

\subsubsection{Actual and warranted growth}

Suppose firms add capacity when they expect sales to expand. Saving responds
to the level of income, while desired investment responds to its expected
increase. Harrod's warranted rate describes growth that validates the resulting
investment plans; his natural rate addresses the scope for sustained growth at
full employment.\footnote{R. F. Harrod (1939), ``An Essay in Dynamic Theory,''
\textit{Economic Journal}, 49(193), 14--33.
\url{https://doi.org/10.2307/2225181}. We distinguish his actual and required
incremental capital coefficients using $v_a$ and $v_r$, and express the argument
in continuous-time notation.}

First distinguish an accounting relation from a behavioral condition. Define
actual growth and, for $\dot Y>0$, an actual incremental capital-output ratio:
\begin{equation}
g_a\equiv\frac{\dot Y}{Y},\qquad v_a\equiv\frac{I}{\dot Y}.
\end{equation}
Because realized saving equals realized investment,
\begin{equation}\label{eq:harrod-actual-identity}
g_av_a=\frac{\dot Y}{Y}\frac{I}{\dot Y}=\frac{I}{Y}=s.
\end{equation}
This identity does not determine $g_a$: $v_a$ changes with the realized outcome.
It is also an incremental ratio, not automatically the stock ratio $K/Y$.

Now let $v_r>0$ be the addition to capital that firms regard as required per
unit of additional output at normal utilization. A simple investment rule is
\begin{equation}
I^d=v_r\dot Y^e,
\end{equation}
where $\dot Y^e$ denotes the expected increase in output. On a path that
validates expectations, $\dot Y^e=\dot Y$ and planned investment equals saving:
\begin{align}
v_r\dot Y&=sY,\\
g_w&\equiv\frac{s}{v_r}.\label{eq:harrod-warranted}
\end{align}
$g_w$ is the warranted growth rate. At this rate, firms' capital additions are
consistent with the realized expansion in sales. Warranted does not mean
socially optimal, forecast with certainty, or necessarily equal to actual growth.

\subsubsection{Why a deviation can amplify}

Compare positive actual growth near $g_w$, holding $s$ and $v_r$ fixed. When
$g_a>g_w$, \eqref{eq:harrod-actual-identity} gives
\begin{equation}
v_a=\frac{s}{g_a}<\frac{s}{g_w}=v_r.
\end{equation}
Capital additions have been small relative to the growth of sales. If firms
interpret this as a shortage of capacity or inventories, they may increase
planned investment. That spending raises demand and can reinforce the expansion.
When $g_a<g_w$, the corresponding $v_a>v_r$ can leave firms with excess capacity
or unwanted inventories; cutting investment then weakens demand further.

The sign of this response is a behavioral assumption, not a consequence of
$S=I$. Realized investment can include unplanned inventory accumulation;
accounting equality does not make all investment plans mutually consistent.
Imposing $I^d=I=sY=v_r\dot Y$ at every date would force $g_a=g_w$ and leave
no off-path adjustment to analyze.

For an explicit \emph{illustrative local} adjustment, set
$x(t)=g_a(t)-g_w$ and assume
\begin{equation}\label{eq:harrod-local-instability}
\dot x=\chi x,\qquad \chi>0
\quad\Longrightarrow\quad x(t)=x(0)e^{\chi t}.
\end{equation}
This formalizes the amplifying response just described. It is not Harrod's
original equation or a global account of the business cycle. Expectations,
financing, utilization, and investment lags could change the adjustment rule.
The identity and warranted-rate condition alone establish neither stability
nor instability.

\subsubsection{Natural growth and the employment constraint}

Let $L$ be the available labor force, growing at $n$. In a simplified
representation, labor productivity at normal operation, $A_L$, grows at
$\gamma$. Sustainable output with full employment then grows at
\begin{equation}\label{eq:harrod-natural}
Y^N=A_LL,\qquad g_n=n+\gamma.
\end{equation}
This is a modern shorthand for the full-employment growth constraint, holding
labor-supply and productivity paths fixed. It is not a limit on every observed
year's GDP growth: an economy can temporarily grow faster by bringing idle
workers into employment.

To see the distinction, assume employed workers $E$ produce $Y=A_LE$.
The employment rate $e=E/L$ then satisfies
\begin{equation}
e=\frac{Y}{A_LL},\qquad \frac{\dot e}{e}=g_a-g_n.
\end{equation}
If warranted growth persistently exceeds natural growth, following the
warranted path would eventually require $e>1$; the path cannot continue under
these labor and productivity assumptions. If warranted growth is below natural
growth, output can follow it while employment fails to keep pace with the labor
force. A path satisfying investment plans need not maintain full employment.

\begin{greenblock}
\textbf{Two different questions.}
The comparison $g_a$ versus $g_w$ concerns adjustment after investment plans are
disappointed. The comparison $g_w$ versus $g_n$ concerns whether a warranted
path can persist with full employment. Their equality does not follow from
accounting, and their difference is not, by itself, a complete instability proof.
\end{greenblock}

\subsection{Fixed proportions, compatibility, and the route to Solow}

A common benchmark collects part of these concerns into a fixed-proportions
production function:
\begin{equation}
Y=\min\left\{\frac{K}{v},\frac{AL}{a}\right\},\qquad v,a>0,
\qquad \frac{\dot A}{A}=\gamma.
\end{equation}
Full use of both inputs requires $K=vY$ and $AL=aY$. Restore depreciation
$\delta\geq0$ and now interpret $I=sY$ as \emph{gross} investment. Differentiating
$K=vY$ and substituting accumulation gives
\begin{align}
v\dot Y&=\dot K=sY-\delta K=(s-\delta v)Y,\\
\frac{\dot Y}{Y}&=\frac{s}{v}-\delta.
\end{align}
The effective-labor requirement also implies $\dot Y/Y=n+\gamma$. Hence
\begin{equation}\label{eq:hd-compatibility}
\frac{s}{v}-\delta=n+\gamma.
\end{equation}
There is no mechanism here that makes independently specified parameters
satisfy this equality. With zero depreciation and compatible investment
coefficients, this combines the warranted and natural rates into $g_w=g_n$.
It is a useful consistency benchmark, not a complete reproduction of both
original papers.

In Domar's simplest constant-productivity case, identifying capacity as
$Q=K/v$ gives $\kappa=1/v$, and the required rate becomes $s\kappa=s/v$.
The matching expression does not erase the difference in questions: Domar
compares demand creation with capacity creation; Harrod also asks how firms
respond when actual growth fails to validate investment plans.

The same algebra $g=s/v$ can even appear in an AK model that assumes all output
is supplied and absorbed. Its presence alone therefore says little about
employment or stability. The conclusions depend on which decisions and
adjustment mechanisms accompany it.

Solow's response is to let factor proportions change, so the capital-output
ratio can adjust rather than remain a fixed recipe. We follow that experiment
in Chapter~\ref{sec:solow}. Maintaining full employment there permits a clear
study of capital adjustment; it does not resolve demand failures by assumption
alone. The present chapter leaves investment behavior and expectations largely
specified from outside. Ramsey will later replace the fixed saving rule with
an intertemporal household choice.
